\section{Introduction}
\label{sec:introduction}

\begin{itemize}

\item Memory performance is critical to many important applications
\item Different applications have different requirements
\item Memory system is getting complicated; so is memory management
\item Coordination between different memory levels is getting important
\item Application-specific memory management is needed
\item Coordinated memory management, including replacement, prefetching
page policy, and scheduling
\item Multi-level coordination
\item Machine-learning based management
\item Feature selection, sharing between different levels, feedback 
between different levels
\item Run-time behavior identification
\item Accelerators for application-specific memory management is needed?
\item What should be the architecture of accelerator?
\item Location: OS, memory controller, DIMM, or chip?

\end{itemize}

Today's computer systems have seen increasing demands on memory systems'
performance, energy efficiency, and reliability~\cite{XXX}. Many 
important applications are memory-intensive, such as XXX~\cite{XXX}.
Innumerous research efforts have been spent to tackle the memory wall
problem~\cite{XXX} at different memory hierarchy levels.  The memory
system is getting more and more complicated, such as more cache levels
to fill the processor-memory speed gap and main memory with diverse
storage technologies.  As a result, the management of memory system
is also getting more and more complicated in order to achieve better
performance, energy efficiency, and reliability.  For instance,
various cache replacement policies have been proposed to reduce the cache miss
rates~\cite{XXX}; numerous prefetching schemes have been proposed to predict 
future cache misses and fetch those in advance~\cite{XXX}; and different memory
access scheduling schemes have been proposed to reorder concurrent main memory
requests in order to reduce their average access latency and increase
their bandwidth~\cite{XXX}.  Although each scheme is effective to
achieve its own goal but multiple schemes usually work in isolation and 
do not consider other schemes employed in the system.  For example,
a replacement policy may not consider the impact on the prefetching
engine or the access scheduler when making the replacement decision.
This is understandable and desirable in some sense since otherwise
its design might get too complicated with high implementation cost.
However, such isolated designs are not globally optimal. To further
advance the memory system performance, coordinated management across
different memory hierarchy levels would be needed.

In this project, we propose a machine learning based, coordinated memory
management framework that coordinates the management policies at different
memory hierarchy levels to achieve the optimal memory performance,
energy efficiency and reliability.  
Our proposed schemes have two steps. The first one
is an enhancement to exsiting systems. It utilizes existing memory
management schemes, such as the existing cache replacement policies or
prefetching policies, but add a helper component at each level to
improve the existing policy's efficiency, such as a prefetching 
filter to filter out some unnecessary prefetching requests.  The helpers
will use the feedback information from other levels to coordinate
the memory management.  This approach can limit the design complexity
since it is an enhancement to existing schemes.  The second approach
will go one step further and re-design the management schemes at
each level by considering the status of other levels.  This would lead
to better performance and energy efficiency but also higher design
complexity and implementation cost.

An important design consideration is how to coordinate the multi-level
cache/memory management schemes and incorporate the feedback
information with the existing memory management schemes, such as
the cache replacement policies and memory scheduling policies.
Different applications or different execution phases of an application
may have different memory access patterns.  As a result, the optimal
cache replacement policy or the optimal memory access scheduling
policy may vary during an application's execution or across different
applications.  The helper component may use the feedback information
from other levels as features of the DNN network to select the
optimal replacement or scheduling policy for the current execution
phase of an application.  In this way, we can utilize the existing
cache/memory management schemes and minimize the implementation
complexity.

We will study which feedback information should be collected at each
level, how to collect them efficiently during run-time, how to pass
them efficiently to other levels, and how to use them to guide the
selection of cache replacement, prefetching and memory access scheduling
policies.

For the second step, we aim to develop a coordinated scheme that optimizes
the management of memory hierarchy globally instead of locally at each 
level.  To achieve this goal, a global controller will collect status
information from each level, determine the features that should be used
at each level, and send the information back to the controller at each
level.  Then each level's local controller will use the selected 
feature to make the decision, such as which cache block to be replaced
or prefetched, and which pending memory request should be served next.
A design challenge is how to balance between the isolation of each level's
management scheme and their coordination.  The isolation helps to reduce
the implementation complexity at each level; while the coordination can
improve the peformance and energy efficiency.  It used to be the isolation
approach dominates since a coordinated management scheme would be too
complicated to be used in practice.  With the wide application of DNN
techniques at architecture designs, the coordination between different
levels with limited extra implementation cost becomes possible.  The
feedback informations from other levels can be used as features of the
management engines at each level.  For instance, the cache replacement
engine can use the feedback information from the memory controller to
aid the selection of replaced blocks.  If one candidate block comes
from an access stream with good locality at DRAM level, while another
candidate block comes from another stream with poor locality, the
cache controller may choose to replace the one with good locality at
DRAM since the former one has smaller miss overhead than the later one.
The memory access scheduler may also use the feedback information from
the cache controller and prefetcher to assist the reordering of 
pending memory requests.  For example, it can give higher priority
to a prefetching request with high confidence level than another
prefetching request with low confidence level.

Of course, all these interactions between different memory hierachy 
levels will increase the complexity of memory management. How to balance
between the design complexity and effectiveness would be an important
part of our study.  Furtheremore, since the complexity of coordinated
memory management would be much more complicated than the conventional
management, we will design accelerators that can speed up the DNN
processing the coordinated memory management.
